\documentclass[11pt,a4paper]{article}

\usepackage[utf8]{inputenc}
\usepackage[T1]{fontenc}
\usepackage[brazilian]{babel}
\usepackage[margin=2cm]{geometry}
\usepackage{enumitem}
\usepackage{hyperref}
\usepackage{titlesec}
\usepackage{parskip}

\hypersetup{
  colorlinks=true,
  linkcolor=black,
  urlcolor=black,
  pdftitle={Kamille Hartmann Maccagnan - Currículo},
  pdfauthor={Kamille Hartmann Maccagnan}
}

\titleformat{\section}{\large\bfseries}{}{0em}{}[\titlerule]
\titlespacing*{\section}{0pt}{12pt}{6pt}

\pagestyle{empty}
\setlength{\parindent}{0pt}
\setlength{\parskip}{6pt}

\newcommand{\cvheader}[4]{%
  \begin{center}
    {\LARGE\bfseries #1}\\[4pt]
    {\large #2}\\[6pt]
    #3\\[2pt]
    #4
  \end{center}
  \vspace{4pt}
}

\newcommand{\experience}[3]{%
  \textbf{#1} \hfill \textit{#2}\\[4pt]
  #3
  \vspace{8pt}
}

\begin{document}

\cvheader{Kamille Hartmann Maccagnan}%
{Business Tech Analyst | Análise de Requisitos | Integração Negócio e Tecnologia}%
{Campo Comprido, Curitiba-PR \textbullet{} (41) 9 9928-4424 \textbullet{} \href{mailto:kamillehartmannm@gmail.com}{kamillehartmannm@gmail.com}}%
{\href{https://linkedin.com/in/kamille-hartmann-maccagnan/}{linkedin.com/in/kamille-hartmann-maccagnan}}

\section*{Resumo Profissional}

Profissional com experiência na interface entre negócio, operações e tecnologia, atuando na análise e documentação de requisitos, investigação de problemas operacionais, padronização de processos e apoio à homologação de entregas. Vivência em ambiente de varejo com múltiplas unidades, com elaboração de documentação funcional, POPs e materiais técnicos para usuários e times de T.I. Experiência em PMO, governança de dados e acompanhamento de projetos, com definição de indicadores, relatórios executivos e monitoramento de performance para visibilidade de entregas e apoio à tomada de decisão. Atuação na organização e detalhamento de demandas, identificação de inconsistências, análise de impactos e facilitação de alinhamentos entre stakeholders de negócio e tecnologia. Perfil analítico, organizado e autônomo, com comunicação clara e inglês avançado.

\section*{Experiência Profissional}

\experience{Anjuss -- Analista de Suporte de TI}{07/2025 -- Atual}{%
Atuação no ecossistema operacional de lojas e sistemas internos, como ponte entre áreas de negócio, unidades de varejo e tecnologia. Responsável pela análise, organização e detalhamento de demandas, com investigação de problemas reportados pelos usuários, identificação de causas operacionais e apoio à resolução de incidentes que impactavam o uso do sistema no dia a dia das lojas. Elaboração de documentação funcional estruturada, incluindo POPs do sistema para as lojas e POPs técnicos para a área de T.I., garantindo clareza de regras, fluxos de uso e critérios de execução para diferentes perfis de usuário. Liderança de treinamentos e onboarding de novos funcionários, com padronização de processos e comunicação assíncrona entre áreas para manter alinhamento e rastreabilidade das solicitações. Desenvolvimento de dashboards em Power BI para monitoramento de indicadores operacionais, apoiando diagnóstico de falhas, visibilidade de performance e priorização de melhorias.}

\experience{HORSE do Brasil -- Estagiária de TI (IS/IT LATAM)}{07/2024 -- 07/2025}{%
Atuação em PMO, governança de dados, suporte e interface com áreas operacionais ligadas ao chão de fábrica, em ambiente multinacional e de alta complexidade. Acompanhamento de portfólio de projetos com organização de cronogramas, indicadores de desempenho e visibilidade de entregas, apoiando a quebra de iniciativas em etapas acompanháveis e na validação de entregas junto a stakeholders de negócio, fornecedores e times técnicos. Elaboração de relatórios executivos e materiais de acompanhamento para apoio à tomada de decisão, com participação em reuniões de alinhamento conduzidas em até três idiomas simultaneamente. Atuação na documentação de processos, análise de riscos, investigação de desvios operacionais e evolução de práticas de gestão, com foco em mitigação de incidentes e melhoria contínua.}

\experience{Restaurante Familiar -- Analista de Dados e Suporte Técnico}{01/2023 -- 07/2024}{%
Atuação na análise de necessidades de gestão do negócio e no diagnóstico de problemas que afetavam a confiabilidade de informações operacionais e financeiras. Mapeamento e estruturação de fluxos de coleta e organização de dados, com criação da base de dados do zero em Excel para viabilizar dashboards em Power BI e definição de indicadores de desempenho. Responsável pela padronização de informações, identificação de inconsistências e investigação de causas que impactavam relatórios e tomada de decisão. Atuação no alinhamento entre áreas operacionais e administrativas, traduzindo demandas do negócio em estruturas de dados e relatórios com critérios claros de validação.}

\section*{Projetos e Conquistas}

Estruturação de KPIs e dashboards executivos para monitoramento de indicadores de desempenho e suporte à tomada de decisão. Participação em iniciativas de melhoria contínua e otimização de processos com foco em eficiência operacional. Automação e padronização de fluxos de dados e relatórios, reduzindo retrabalho e aumentando a confiabilidade das informações. 1º lugar no Transformation Day Renault 2023 com solução orientada a dados e melhoria de processos. Participação no Hackathon Bosch 2023 com foco em resolução de problemas de negócio com uso de dados.

\section*{Habilidades Técnicas}

\textbf{Análise Técnica de Negócios e Requisitos:} Levantamento e detalhamento de demandas, documentação funcional, elaboração de POPs e manuais, mapeamento de processos, análise de impactos, validação de entregas, suporte à homologação funcional, rastreabilidade de solicitações.

\textbf{Diagnóstico e Operação:} Investigação de problemas operacionais, identificação de inconsistências, análise de causas, monitoramento de indicadores, apoio à mitigação de incidentes, melhoria contínua.

\textbf{Indicadores e Visibilidade:} KPIs, dashboards executivos, relatórios gerenciais, Power BI (DAX, Power Query), Excel, SQL, governança de dados.

\textbf{Gestão de Demandas e Projetos:} PMO, organização e priorização de demandas, acompanhamento de entregas, Scrum, Kanban, Boas Práticas PMOBOK.

\textbf{Ferramentas:} Monday, Trello, ServiceNow, Spotfire, Pacote Office.

\textbf{Comunicação e Atendimento:} Interface entre negócio e tecnologia, comunicação com stakeholders, suporte funcional, treinamento e onboarding, inglês avançado.

\textbf{Idiomas:} Inglês avançado \textbullet{} Espanhol básico \textbullet{} Coreano básico.

\section*{Formação Acadêmica}

\textbf{PUCPR} -- Graduação em Gestão da Tecnologia da Informação \hfill Previsão de conclusão: 06/2026\\
Curitiba, PR

\textbf{Escola Conquer} -- Pós-graduação em Liderança e Gestão em Tecnologia \hfill Conclusão: 2024\\
Curitiba, PR

\textbf{Universidade Tuiuti do Paraná} -- Graduação em Medicina Veterinária \hfill Conclusão: 06/2019\\
Curitiba, PR

\textbf{Qualittas} -- Pós-graduação em Clínica Médica e Cirúrgica de Animais Exóticos e Pets Não Convencionais \hfill Conclusão: 12/2022\\
Curitiba, PR

\end{document}
